% ch4-determinants.tex — Chapter 4: Determinants
% Permutations, Leibniz formula, cofactor expansion, adjugate, Cramer's rule,
% geometric interpretation, applications.

\chapter{Determinants}\label{ch:determinants}

% ============================================================
% Motivating introduction
% ============================================================

How can one tell, by a single number, whether a square matrix is
invertible? How can one compute the area of a parallelogram Spanned by
two vectors, or the volume of a parallelepiped in~$\R^3$? These
seemingly different questions are all answered by the same object:
the \emph{determinant}\index{determinant}.

Consider two vectors $\vec{u}=(a,b)$ and $\vec{v}=(c,d)$ in~$\R^2$.
The signed area of the parallelogram they Span is
$ad - bc$---precisely the determinant of the matrix whose rows (or
columns) are $\vec{u}$ and~$\vec{v}$. This quantity is zero exactly
when the two vectors are collinear, i.e.\ when the matrix
$\begin{psmallmatrix} a & b \\ c & d \end{psmallmatrix}$ fails to be
invertible. The sign records the \emph{orientation}: positive if
$\vec{v}$ lies to the left of~$\vec{u}$, negative otherwise.

\begin{center}
\begin{tikzpicture}[>=Stealth, scale=1.1]
  % Parallelogram
  \coordinate (O) at (0,0);
  \coordinate (U) at (3,0.8);
  \coordinate (V) at (1,2.2);
  \coordinate (UV) at ($(U)+(V)$);

  \fill[blue!10] (O) -- (U) -- (UV) -- (V) -- cycle;
  \draw[thick, blue!70!black, ->] (O) -- (U) node[midway, below right] {$\vec{u}$};
  \draw[thick, red!70!black, ->] (O) -- (V) node[midway, above left] {$\vec{v}$};
  \draw[dashed, gray] (U) -- (UV) -- (V);

  \node at ($(O)!0.5!(UV)$) {$\abs{\det}$};

  % Orientation arc
  \draw[->, thick, green!50!black] (0.9,0.24) arc[start angle=15, end angle=55, radius=0.9]
    node[midway, right] {\small $+$};
\end{tikzpicture}
\captionof{figure}{The absolute value of $\det\begin{psmallmatrix}a&b\\c&d\end{psmallmatrix}$
  gives the area of the parallelogram; the sign encodes orientation.}
\label{fig:det-parallelogram-intro}
\end{center}

In this chapter we develop the theory of determinants rigorously. We
begin with the combinatorial machinery of permutations and their
signatures, define the determinant via the Leibniz formula, establish
its fundamental properties, prove cofactor expansion, derive Cramer's
rule, and explore geometric interpretations and applications.

Throughout this chapter, $\K$ denotes a field (typically $\R$ or~$\C$),
and all matrices are square unless stated otherwise. We write
$\Mat_n(\K)$ for the space of $n\times n$ matrices with entries
in~$\K$, and we freely use notation and results from

ef{ch:vector-spaces,ch:linear-maps}.


% ============================================================
\section{Permutations and the symmetric group}\label{sec:permutations}
% ============================================================

\begin{definition}{Permutation}{def:permutation}
\index{permutation}\index{symmetric group}
A \emph{permutation} of the set $\set{1,2,\ldots,n}$ is a bijection
$\sigma\colon\set{1,\ldots,n}\to\set{1,\ldots,n}$.
The set of all such permutations, equipped with composition, is the
\emph{symmetric group}\index{symmetric group} $\mathfrak{S}_n$.
Its cardinality is $\card{\mathfrak{S}_n} = n!$.
\end{definition}

We use two-line notation and cycle notation interchangeably. For
instance, the permutation $\sigma\in\mathfrak{S}_4$ with
$\sigma(1)=2$, $\sigma(2)=4$, $\sigma(3)=3$, $\sigma(4)=1$ is written
\[
  \sigma = \begin{pmatrix} 1 & 2 & 3 & 4 \\ 2 & 4 & 3 & 1 \end{pmatrix}
  = (1\;2\;4).
\]

\begin{definition}{Transposition}{def:transposition}
\index{transposition}
A \emph{transposition} is a permutation that exchanges exactly two
elements and fixes all others. We write $\tau_{ij}$ for the
transposition that swaps $i$ and~$j$ (with $i\neq j$).
\end{definition}

\begin{proposition}{Generation by transpositions}{prop:gen-transpositions}
Every permutation $\sigma\in\mathfrak{S}_n$ can be written as a
product (composition) of transpositions.
\end{proposition}

\begin{proof}
It suffices to show every cycle decomposes into transpositions, since
every permutation is a product of disjoint cycles. Indeed,
$(a_1\;a_2\;\cdots\;a_k) = \tau_{a_1 a_k}\circ\tau_{a_1 a_{k-1}}
\circ\cdots\circ\tau_{a_1 a_2}$.
\end{proof}

\begin{definition}{Inversions and signature}{def:signature}
\index{inversion}\index{signature!of a permutation}
An \emph{inversion} of $\sigma\in\mathfrak{S}_n$ is a pair $(i,j)$
with $1\leq i < j \leq n$ and $\sigma(i)>\sigma(j)$.
The \emph{signature} (or \emph{sign}) of~$\sigma$ is
\[
  \sgn(\sigma) \deff (-1)^{N(\sigma)},
\]
where $N(\sigma)$ denotes the number of inversions of~$\sigma$.
Equivalently,
\[
  \sgn(\sigma) = \prod_{1\leq i < j \leq n}
    \frac{\sigma(j)-\sigma(i)}{j-i}.
\]
A permutation is called \emph{even}\index{even permutation} if
$\sgn(\sigma)=+1$ and \emph{odd}\index{odd permutation} if
$\sgn(\sigma)=-1$.
\end{definition}

\begin{proposition}{Properties of the signature}{prop:sgn-properties}
For all $\sigma,\tau\in\mathfrak{S}_n$:
\begin{enumerate}[label=(\roman*)]
  \item $\sgn(\sigma\circ\tau) = \sgn(\sigma)\cdot\sgn(\tau)$.
  \item $\sgn(\sigma^{-1}) = \sgn(\sigma)$.
  \item Every transposition has signature $-1$.
  \item If $\sigma = \tau_1\circ\cdots\circ\tau_r$ is a product of
    transpositions, then $\sgn(\sigma)=(-1)^r$. In particular, the
    parity of~$r$ depends only on~$\sigma$, not on the decomposition.
\end{enumerate}
\end{proposition}

\begin{proof}
Consider the product formula
$\sgn(\sigma) = \prod_{i<j}\frac{\sigma(j)-\sigma(i)}{j-i}$.

(iii) Let $\tau = \tau_{pq}$ with $p < q$. In the product
$\prod_{i<j}(\tau(j)-\tau(i))$, compared to $\prod_{i<j}(j-i)$,
the factor $(q-p)$ changes sign (to $(p-q)$), and for each $k$
strictly between $p$ and $q$, the two factors involving $k$ and $p$,
and $k$ and $q$, exchange roles, contributing no net sign change.
A careful count shows the total sign is $-1$.

(i) From the product formula,
$\sgn(\sigma\circ\tau) = \prod_{i<j}
\frac{(\sigma\circ\tau)(j)-(\sigma\circ\tau)(i)}{j-i}$.
Write $j-i = \frac{j-i}{\tau(j)-\tau(i)}\cdot(\tau(j)-\tau(i))$
(since $\tau$ is a bijection, the denominator is nonzero when $i<j$),
and rearrange to factor the product as $\sgn(\sigma)\cdot\sgn(\tau)$.

(ii) Follows from (i) applied to $\sigma\circ\sigma^{-1}=\Id$
and $\sgn(\Id)=+1$.

(iv) Follows from (i) and (iii) by induction.
\end{proof}

\begin{example}{Signatures in $\mathfrak{S}_3$}{ex:S3-signatures}
The group $\mathfrak{S}_3$ has $3!=6$ elements:
\[
\begin{array}{c|cccccc}
  \sigma & \Id & (1\;2) & (1\;3) & (2\;3) & (1\;2\;3) & (1\;3\;2) \\
  \hline
  \sgn(\sigma) & +1 & -1 & -1 & -1 & +1 & +1
\end{array}
\]
The three-cycles are even (each is a product of two transpositions:
$(1\;2\;3)=(1\;3)(1\;2)$), while the transpositions are odd.
\end{example}

\begin{example}{Computing a signature}{ex:signature-computation}
Let $\sigma = (1\;3\;5\;2)\in\mathfrak{S}_5$, so $\sigma(1)=3$,
$\sigma(3)=5$, $\sigma(5)=2$, $\sigma(2)=1$, $\sigma(4)=4$.
The inversions are $(1,2),(2,5),(3,5)$, giving $N(\sigma)=3$.
Alternatively, $(1\;3\;5\;2)=(1\;2)(1\;5)(1\;3)$ is a product of
3~transpositions, so $\sgn(\sigma)=(-1)^3 = -1$.
\end{example}


% ============================================================
\section{Definition of the determinant}\label{sec:det-definition}
% ============================================================

\begin{definition}{Determinant (Leibniz formula)}{def:determinant}
\index{determinant!Leibniz formula}\index{Leibniz formula}
Let $A = (a_{ij})_{1\leq i,j\leq n}\in\Mat_n(\K)$. The
\emph{determinant} of~$A$ is
\[
  \det(A) \deff \sum_{\sigma\in\mathfrak{S}_n}
    \sgn(\sigma)\,\prod_{i=1}^n a_{i,\sigma(i)}.
\]
\end{definition}

\begin{remark}{Small cases}{rem:small-det}
For $n=1$: $\det(a) = a$.

For $n=2$:
$\det\begin{pmatrix} a & b \\ c & d \end{pmatrix} = ad - bc$.

For $n=3$:
$\det\begin{pmatrix} a_1 & a_2 & a_3 \\ b_1 & b_2 & b_3 \\
c_1 & c_2 & c_3 \end{pmatrix}
= a_1 b_2 c_3 + a_2 b_3 c_1 + a_3 b_1 c_2
  - a_3 b_2 c_1 - a_2 b_1 c_3 - a_1 b_3 c_2$,
which is the \emph{Sarrus rule}\index{Sarrus rule}.
\end{remark}


% ============================================================
\section{Alternating multilinear characterization}\label{sec:multilinear}
% ============================================================

We now give an axiomatic characterization of the determinant. Denote
the rows of $A$ by $R_1,\ldots,R_n\in\K^n$, so that we may write
$\det(A)$ as a function of these row vectors:
$\det(A) = D(R_1,\ldots,R_n)$.

\begin{definition}{Multilinear and alternating}{def:multilinear-alternating}
\index{multilinear form}\index{alternating form}
A function $D\colon (\K^n)^n \to \K$ is:
\begin{enumerate}[label=(\roman*)]
  \item \emph{multilinear} (in the rows) if for each $i$,
    $D$ is linear in the $i$-th argument with all other arguments fixed:
    \[
      D(R_1,\ldots,\alpha R_i + \beta R_i',\ldots,R_n)
      = \alpha\,D(R_1,\ldots,R_i,\ldots,R_n)
        + \beta\,D(R_1,\ldots,R_i',\ldots,R_n).
    \]
  \item \emph{alternating} if $D(R_1,\ldots,R_n)=0$ whenever two rows
    are equal: $R_i = R_j$ for some $i\neq j$.
\end{enumerate}
\end{definition}

\begin{lemma}{Alternating implies antisymmetric}{lem:alt-antisymmetric}
If $D$ is multilinear and alternating, then swapping two rows changes
the sign:
\[
  D(R_1,\ldots,R_i,\ldots,R_j,\ldots,R_n)
  = -D(R_1,\ldots,R_j,\ldots,R_i,\ldots,R_n).
\]
\end{lemma}

\begin{proof}
Expand $D(\ldots,R_i+R_j,\ldots,R_i+R_j,\ldots) = 0$ by multilinearity.
Two of the four terms vanish (identical rows), leaving
$D(\ldots,R_i,\ldots,R_j,\ldots) + D(\ldots,R_j,\ldots,R_i,\ldots) = 0$.
\end{proof}

\begin{theorem}{Uniqueness of the determinant}{thm:det-unique}
\index{determinant!axiomatic characterization}
The determinant is the unique function
$D\colon (\K^n)^n \to \K$ that is:
\begin{enumerate}[label=(\roman*)]
  \item multilinear in the rows,
  \item alternating,
  \item normalized: $D(I_n)=1$.
\end{enumerate}
\end{theorem}

\begin{proof}
\textbf{Existence.} We verify that the Leibniz formula defines a
function satisfying (i)--(iii).

Multilinearity: each term
$\sgn(\sigma)\prod_{i=1}^n a_{i,\sigma(i)}$ is linear in each
row~$R_k$ (the factor $a_{k,\sigma(k)}$ appears linearly), and the
sum of linear functions is linear.

Alternating: suppose $R_p = R_q$ for $p\neq q$. Pair each
$\sigma\in\mathfrak{S}_n$ with $\sigma'=\sigma\circ\tau_{pq}$.
Then $\sgn(\sigma')=-\sgn(\sigma)$ while
$\prod_i a_{i,\sigma'(i)}=\prod_i a_{i,\sigma(i)}$ (since
swapping rows $p,q$ with identical entries has no effect on the
product). Hence the terms cancel pairwise, giving $D=0$.

Normalization: for $A=I_n$, $a_{i,\sigma(i)}=\delta_{i,\sigma(i)}$,
so the only surviving permutation is $\sigma=\Id$, giving
$D(I_n)=\sgn(\Id)\cdot 1 = 1$.

\textbf{Uniqueness.} Let $D$ satisfy (i)--(iii). Write
$R_i = \sum_{j=1}^n a_{ij}\,\vece{j}$. By multilinearity,
\[
  D(R_1,\ldots,R_n)
  = \sum_{j_1,\ldots,j_n=1}^{n}
    a_{1,j_1}\cdots a_{n,j_n}\,D(\vece{j_1},\ldots,\vece{j_n}).
\]
By the alternating property, $D(\vece{j_1},\ldots,\vece{j_n})=0$
unless $j_1,\ldots,j_n$ are all distinct, i.e.\ $(j_1,\ldots,j_n)$
is a permutation $\sigma$. In that case, repeated application of

ef{lem:lem:alt-antisymmetric} gives
$D(\vece{\sigma(1)},\ldots,\vece{\sigma(n)})
= \sgn(\sigma)\,D(\vece{1},\ldots,\vece{n})
= \sgn(\sigma)$.
Therefore $D(R_1,\ldots,R_n) = \sum_{\sigma\in\mathfrak{S}_n}
\sgn(\sigma)\prod_{i=1}^n a_{i,\sigma(i)}$, which is the Leibniz
formula.
\end{proof}


% ============================================================
\section{Properties of the determinant}\label{sec:det-properties}
% ============================================================

\begin{proposition}{Effect of row operations}{prop:row-ops}
\index{determinant!row operations}
Let $A\in\Mat_n(\K)$ with rows $R_1,\ldots,R_n$.
\begin{enumerate}[label=(\roman*)]
  \item \textbf{Scaling:} Multiplying row $R_i$ by $\lambda\in\K$
    multiplies $\det(A)$ by $\lambda$.
  \item \textbf{Swap:} Swapping two rows multiplies $\det(A)$ by $-1$.
  \item \textbf{Row addition:} Adding a scalar multiple of one row to
    another does not change $\det(A)$:
    replacing $R_i$ by $R_i + \lambda R_j$ (with $j\neq i$) leaves
    $\det$ unchanged.
\end{enumerate}
In particular, $\det(\lambda A)=\lambda^n\det(A)$ for all
$\lambda\in\K$.
\end{proposition}

\begin{proof}
(i) and (ii) follow immediately from multilinearity and

ef{lem:lem:alt-antisymmetric}.

(iii) By multilinearity in row~$i$,
$D(R_1,\ldots,R_i+\lambda R_j,\ldots,R_n)
= D(R_1,\ldots,R_i,\ldots,R_n)
+ \lambda\,D(R_1,\ldots,R_j,\ldots,R_n)$,
where in the second term row~$j$ appears in both positions $i$
and~$j$, so the alternating property gives $0$.
\end{proof}

\begin{theorem}{Determinant of a transpose}{thm:det-transpose}
\index{determinant!of transpose}
For all $A\in\Mat_n(\K)$,
$\det(\transpose{A}) = \det(A)$.
\end{theorem}

\begin{proof}
Denote $B = \transpose{A}$, so $b_{ij}=a_{ji}$. Then
\begin{align*}
  \det(B) &= \sum_{\sigma\in\mathfrak{S}_n}\sgn(\sigma)
    \prod_{i=1}^n b_{i,\sigma(i)}
  = \sum_{\sigma\in\mathfrak{S}_n}\sgn(\sigma)
    \prod_{i=1}^n a_{\sigma(i),i}.
\end{align*}
Substituting $\tau = \sigma^{-1}$ (a bijection on $\mathfrak{S}_n$
with $\sgn(\tau)=\sgn(\sigma)$), and writing $j=\sigma(i)$
so $i=\tau(j)$:
\[
  \det(B) = \sum_{\tau\in\mathfrak{S}_n}\sgn(\tau)
    \prod_{j=1}^n a_{j,\tau(j)} = \det(A). \qedhere
\]
\end{proof}

\begin{remark}{Columns behave like rows}{rem:columns}
Since $\det(\transpose{A})=\det(A)$, every property of $\det$ stated
for rows also holds for columns: the determinant is also multilinear
and alternating in the columns.
\end{remark}

\begin{theorem}{Determinant of a product}{thm:det-product}
\index{determinant!of product}
For all $A,B\in\Mat_n(\K)$,
\[
  \det(AB) = \det(A)\det(B).
\]
\end{theorem}

\begin{proof}
Fix $A\in\Mat_n(\K)$ and define $D\colon\Mat_n(\K)\to\K$ by
$D(B) = \det(AB)$, viewed as a function of the rows of~$B$.

If $B$ has rows $R_1,\ldots,R_n$, then $AB$ has rows
$AR_1^{\mathsf{T}},\ldots,AR_n^{\mathsf{T}}$ (where we view each
$R_i$ as a row vector and $A$ acts by left-multiplication after
transposing). More precisely, the $i$-th row of $AB$ is
$\sum_{k=1}^n b_{ik}\,C_k$, where $C_k$ is the $k$-th row of~$A$
viewed appropriately. Rewriting: the $i$-th row of $AB$ is a linear
function of the $i$-th row of~$B$.

Hence $B\mapsto\det(AB)$ is multilinear and alternating in the
rows of~$B$ (since $\det$ itself is). By the uniqueness

ef{thm:thm:det-unique},
$D(B) = D(I_n)\cdot\det(B) = \det(A)\cdot\det(B)$.
\end{proof}

\begin{corollary}{Determinant of an inverse}{cor:det-inverse}
\index{determinant!of inverse}
If $A\in\GL_n(\K)$, then $\det(\inv{A}) = \frac{1}{\det(A)}$.
\end{corollary}

\begin{proof}
$\det(A)\det(\inv{A}) = \det(A\inv{A}) = \det(I_n)=1$.
\end{proof}

\begin{proposition}{Triangular matrices}{prop:det-triangular}
\index{determinant!triangular matrix}\index{triangular matrix!determinant}
If $A$ is upper or lower triangular, then
$\det(A) = \prod_{i=1}^n a_{ii}$.
In particular, $\det(I_n)=1$ and
$\det(\diag(\lambda_1,\ldots,\lambda_n)) = \lambda_1\cdots\lambda_n$.
\end{proposition}

\begin{proof}
In the Leibniz formula, the product $\prod_{i=1}^n a_{i,\sigma(i)}$
vanishes unless $\sigma(i)\geq i$ for all~$i$ (upper triangular case).
The only permutation satisfying this is $\sigma=\Id$, so
$\det(A) = \prod_{i=1}^n a_{ii}$.
\end{proof}

\begin{proposition}{Block triangular matrices}{prop:det-block}
\index{determinant!block matrix}
Let $A$ be a block upper (or lower) triangular matrix:
\[
  A = \begin{pmatrix} A_{11} & A_{12} \\ 0 & A_{22} \end{pmatrix},
\]
where $A_{11}\in\Mat_p(\K)$ and $A_{22}\in\Mat_q(\K)$ with $p+q=n$.
Then $\det(A) = \det(A_{11})\cdot\det(A_{22})$.
\end{proposition}

\begin{proof}
Fix $A_{11}$ and $A_{12}$, and consider
$D(A_{22}) \deff \det(A)$ as a function of the rows of $A_{22}$
(these are the last $q$ rows of $A$). Since the block has zeros below
$A_{11}$, the last $q$ rows of $A$ are
$(0,\ldots,0,r_{q,1},\ldots,r_{q,q})$ where $(r_{q,1},\ldots,r_{q,q})$
is a row of~$A_{22}$. By multilinearity and the alternating property
of $\det$, $D$ is multilinear and alternating in the rows of $A_{22}$.
Hence $D(A_{22}) = D(I_q)\cdot\det(A_{22})$. But when $A_{22}=I_q$,
the matrix $A$ is block upper triangular with $A_{22}=I_q$, and we
can row-reduce $A_{12}$ to zero using the last $q$ rows without
changing the determinant, obtaining
$\det\begin{psmallmatrix} A_{11} & 0 \\ 0 & I_q\end{psmallmatrix}
= \det(A_{11})$.
Therefore $D(I_q)=\det(A_{11})$.
\end{proof}


% ============================================================
\section{Cofactor expansion (Laplace expansion)}\label{sec:cofactor}
% ============================================================

\begin{definition}{Minor and cofactor}{def:cofactor}
\index{minor}\index{cofactor}
Let $A=(a_{ij})\in\Mat_n(\K)$. The \emph{$(i,j)$-minor} of~$A$,
denoted $M_{ij}$, is the determinant of the $(n-1)\times(n-1)$
submatrix obtained by deleting row~$i$ and column~$j$. The
\emph{$(i,j)$-cofactor} is
\[
  C_{ij} \deff (-1)^{i+j}\,M_{ij}.
\]
\end{definition}

\begin{theorem}{Laplace expansion}{thm:laplace}
\index{Laplace expansion}\index{cofactor expansion}
For any $A\in\Mat_n(\K)$:
\begin{enumerate}[label=(\roman*)]
  \item \textbf{Expansion along row~$i$:}\quad
    $\displaystyle\det(A) = \sum_{j=1}^n a_{ij}\,C_{ij}
     = \sum_{j=1}^n (-1)^{i+j}\,a_{ij}\,M_{ij}$.
  \item \textbf{Expansion along column~$j$:}\quad
    $\displaystyle\det(A) = \sum_{i=1}^n a_{ij}\,C_{ij}
     = \sum_{i=1}^n (-1)^{i+j}\,a_{ij}\,M_{ij}$.
\end{enumerate}
Moreover, expanding along a ``foreign'' row or column gives zero:
\[
  \sum_{j=1}^n a_{ij}\,C_{kj} = 0 \quad\text{for } i\neq k,
  \qquad
  \sum_{i=1}^n a_{ij}\,C_{ik} = 0 \quad\text{for } j\neq k.
\]
\end{theorem}

\begin{proof}
\textbf{(i) Expansion along row~$i$.}
In the Leibniz formula $\det(A)=\sum_{\sigma}\sgn(\sigma)
\prod_{\ell=1}^n a_{\ell,\sigma(\ell)}$, group the terms according
to the value $\sigma(i)=j$:
\[
  \det(A) = \sum_{j=1}^n a_{ij}
    \sum_{\substack{\sigma\in\mathfrak{S}_n\\\sigma(i)=j}}
    \sgn(\sigma)\prod_{\ell\neq i}a_{\ell,\sigma(\ell)}.
\]
We claim that
$\sum_{\sigma:\sigma(i)=j}\sgn(\sigma)\prod_{\ell\neq i}
a_{\ell,\sigma(\ell)} = (-1)^{i+j}M_{ij}$.

To see this, consider the permutations $\sigma\in\mathfrak{S}_n$
with $\sigma(i)=j$. Any such $\sigma$ is determined by its restriction
to $\set{1,\ldots,n}\setminus\set{i}$, which is a bijection onto
$\set{1,\ldots,n}\setminus\set{j}$. We can identify this restriction
with a permutation $\sigma'\in\mathfrak{S}_{n-1}$ after relabelling:
the rows $1,\ldots,i-1,i+1,\ldots,n$ become $1,\ldots,n-1$ and
similarly for columns. The relabelling involves moving index~$i$ to
position~$n$ (via $i-1$ adjacent transpositions in the row indices)
and index~$j$ to position~$n$ (via $j-1$ adjacent transpositions
in the column indices), contributing a sign of
$(-1)^{(n-i)+(n-j)}=(-1)^{i+j}$ (since $(-1)^{2n}=1$).
Thus $\sgn(\sigma) = (-1)^{i+j}\sgn(\sigma')$, and the inner sum
becomes $(-1)^{i+j}\det(A_{ij})=(-1)^{i+j}M_{ij}=C_{ij}$.

\textbf{(ii)} follows from (i) applied to $\transpose{A}$, using
$\det(\transpose{A})=\det(A)$.

\textbf{Foreign expansion.}
Consider $\sum_{j=1}^n a_{ij}C_{kj}$ with $i\neq k$. This equals the
determinant of the matrix $A'$ obtained from $A$ by replacing row~$k$
with row~$R_i$ (since the cofactors $C_{kj}$ depend only on the rows
other than row~$k$, and the formula gives the expansion of $\det(A')$
along row~$k$). But $A'$ has two identical rows ($R_i$ appears in both
positions $i$ and~$k$), so $\det(A')=0$ by the alternating property.
\end{proof}

\begin{example}{$3\times 3$ determinant by cofactor expansion}{ex:3x3-cofactor}
Compute $\det\begin{pmatrix} 2 & 1 & 3 \\ 0 & -1 & 4 \\ 5 & 2 & 1
\end{pmatrix}$ by expanding along the first column:
\begin{align*}
  \det(A) &= 2\cdot(-1)^{1+1}\det\begin{pmatrix}-1&4\\2&1\end{pmatrix}
  + 0\cdot(-1)^{2+1}\det\begin{pmatrix}1&3\\2&1\end{pmatrix}
  + 5\cdot(-1)^{3+1}\det\begin{pmatrix}1&3\\-1&4\end{pmatrix}\\
  &= 2(-1-8) + 0 + 5(4+3) = 2(-9)+5\cdot 7 = -18+35 = 17.
\end{align*}
\end{example}


% ============================================================
\section{The adjugate matrix and Cramer's rule}\label{sec:adjugate}
% ============================================================

\begin{definition}{Adjugate (classical adjoint)}{def:adjugate}
\index{adjugate matrix}\index{comatrix}
The \emph{adjugate} (or \emph{classical adjoint}, or
\emph{comatrix}\index{comatrix|see{adjugate matrix}}) of
$A\in\Mat_n(\K)$ is the $n\times n$ matrix
\[
  \operatorname{adj}(A) \deff \bigl(C_{ji}\bigr)_{1\leq i,j\leq n}
  = \transpose{(\text{cofactor matrix})},
\]
where $C_{ij}=(-1)^{i+j}M_{ij}$ is the $(i,j)$-cofactor.
\end{definition}

\begin{theorem}{Adjugate identity}{thm:adjugate}
\index{adjugate matrix!identity}
For all $A\in\Mat_n(\K)$,
\[
  A\cdot\operatorname{adj}(A) = \operatorname{adj}(A)\cdot A
  = \det(A)\cdot I_n.
\]
\end{theorem}

\begin{proof}
The $(i,k)$-entry of $A\cdot\operatorname{adj}(A)$ is
$\sum_{j=1}^n a_{ij}\,[\operatorname{adj}(A)]_{jk}
= \sum_{j=1}^n a_{ij}\,C_{kj}$.
By 
ef{thm:thm:laplace}, this equals $\det(A)$ if $i=k$ and $0$ if
$i\neq k$ (foreign expansion). Thus
$A\cdot\operatorname{adj}(A)=\det(A)\,I_n$. The identity
$\operatorname{adj}(A)\cdot A = \det(A)\,I_n$ is proved by
considering columns instead of rows (or by applying the row result
to $\transpose{A}$ and transposing).
\end{proof}

\begin{corollary}{Inverse via adjugate}{cor:inverse-adjugate}
\index{inverse matrix!via adjugate}
If $\det(A)\neq 0$, then $A$ is invertible and
\[
  \inv{A} = \frac{1}{\det(A)}\operatorname{adj}(A).
\]
\end{corollary}

\begin{theorem}{Cramer's rule}{thm:cramer}
\index{Cramer's rule}
Let $A\in\Mat_n(\K)$ with $\det(A)\neq 0$, and let $b\in\K^n$.
The unique solution $x=(x_1,\ldots,x_n)^{\mathsf{T}}$ of $Ax=b$ is
given by
\[
  x_j = \frac{\det(A_j)}{\det(A)},\qquad j=1,\ldots,n,
\]
where $A_j$ is the matrix obtained from~$A$ by replacing column~$j$
with~$b$.
\end{theorem}

\begin{proof}
Since $\det(A)\neq 0$, the system has a unique solution
$x = \inv{A}b = \frac{1}{\det(A)}\operatorname{adj}(A)\,b$.
The $j$-th component is
\[
  x_j = \frac{1}{\det(A)}\sum_{i=1}^n C_{ij}\,b_i
  = \frac{1}{\det(A)}\sum_{i=1}^n (-1)^{i+j}M_{ij}\,b_i.
\]
But $\sum_{i=1}^n (-1)^{i+j}b_i\,M_{ij}$ is precisely the cofactor
expansion of $\det(A_j)$ along column~$j$ (since column~$j$ of $A_j$
is~$b$, and the minors $M_{ij}$ are unchanged). Hence
$x_j = \det(A_j)/\det(A)$.
\end{proof}

\begin{example}{Cramer's rule in dimension 2}{ex:cramer-2d}
Solve $\begin{cases} 3x+2y=7 \\ x-y=1\end{cases}$.

Here $A=\begin{psmallmatrix}3&2\\1&-1\end{psmallmatrix}$,
$\det(A)=-3-2=-5$,
$\det(A_1)=\det\begin{psmallmatrix}7&2\\1&-1\end{psmallmatrix}=-9$,
$\det(A_2)=\det\begin{psmallmatrix}3&7\\1&1\end{psmallmatrix}=-4$.
So $x = \frac{-9}{-5}=\frac{9}{5}$,
$y=\frac{-4}{-5}=\frac{4}{5}$.
\end{example}


% ============================================================
\section{Determinant and invertibility}\label{sec:det-invertibility}
% ============================================================

\begin{theorem}{Invertibility criterion}{thm:det-invertibility}
\index{determinant!and invertibility}\index{invertible matrix!determinant criterion}
A matrix $A\in\Mat_n(\K)$ is invertible if and only if $\det(A)\neq 0$.
\end{theorem}

\begin{proof}
$(\Rightarrow)$\; If $A$ is invertible, then
$\det(A)\det(\inv{A})=\det(I_n)=1$, so $\det(A)\neq 0$.

$(\Leftarrow)$\; If $\det(A)\neq 0$, then by 
ef{thm:thm:adjugate},
$A\cdot\frac{1}{\det(A)}\operatorname{adj}(A) = I_n$,
so $A$ is invertible (with explicit inverse given by

ef{cor:cor:inverse-adjugate}).

Alternatively, if $A$ is not invertible, its columns are linearly
dependent: some column is a linear combination of the others. By
multilinearity and the alternating property in columns,
$\det(A)=0$.
\end{proof}

\begin{corollary}{The general linear group}{cor:GL-det}
\index{general linear group}
$\GL_n(\K) = \setcond{A\in\Mat_n(\K)}{\det(A)\neq 0}$.
The determinant defines a group homomorphism
$\det\colon\GL_n(\K)\to\K^*$ (where $\K^*=\K\setminus\set{0}$
is the multiplicative group of~$\K$).
\end{corollary}

\begin{proof}
The first statement is 
ef{thm:thm:det-invertibility}. The
multiplicativity $\det(AB)=\det(A)\det(B)$
(
ef{thm:thm:det-product}) shows that $\det$ is a group homomorphism.
Its kernel is the \emph{special linear group}
$\SL_n(\K)=\setcond{A\in\GL_n(\K)}{\det(A)=1}$.
\end{proof}


% ============================================================
\section{Determinant of a linear map}\label{sec:det-linear-map}
% ============================================================

\begin{definition}{Determinant of an endomorphism}{def:det-endo}
\index{determinant!of a linear map}
Let $E$ be a finite-dimensional $\K$-vector space and
$f\in\End(E)$. Choose any basis $\cB$ of~$E$ and define
\[
  \det(f) \deff \det\bigl(\Mat_\cB(f)\bigr).
\]
\end{definition}

\begin{proposition}{Well-definedness}{prop:det-well-defined}
The determinant $\det(f)$ does not depend on the choice of basis.
\end{proposition}

\begin{proof}
Let $\cB$ and $\cB'$ be two bases of~$E$, and let $P$ be the
change-of-basis matrix from $\cB$ to $\cB'$. Then
$\Mat_{\cB'}(f) = \inv{P}\,\Mat_\cB(f)\,P$, so
\[
  \det\bigl(\Mat_{\cB'}(f)\bigr)
  = \det(\inv{P})\det\bigl(\Mat_\cB(f)\bigr)\det(P)
  = \frac{1}{\det(P)}\det\bigl(\Mat_\cB(f)\bigr)\det(P)
  = \det\bigl(\Mat_\cB(f)\bigr). \qedhere
\]
\end{proof}

\begin{proposition}{Properties of $\det$ for endomorphisms}{prop:det-endo-props}
Let $E$ be a finite-dimensional $\K$-vector space and
$f,g\in\End(E)$.
\begin{enumerate}[label=(\roman*)]
  \item $\det(g\circ f) = \det(g)\det(f)$.
  \item $f$ is an automorphism if and only if $\det(f)\neq 0$.
  \item $\det(\Id_E)=1$.
  \item $\det(\lambda f)=\lambda^n\det(f)$ where $n=\dim E$.
\end{enumerate}
\end{proposition}

\begin{proof}
All properties follow immediately from the corresponding properties
of matrix determinants and the identity
$\Mat_\cB(g\circ f)=\Mat_\cB(g)\cdot\Mat_\cB(f)$.
\end{proof}


% ============================================================
\section{Geometric interpretation}\label{sec:det-geometry}
% ============================================================

The determinant has a rich geometric meaning in~$\R^n$.

\begin{theorem}{Determinant as signed volume}{thm:signed-volume}
\index{determinant!geometric interpretation}\index{signed volume}
Let $v_1,\ldots,v_n\in\R^n$ be the column vectors of a matrix
$A\in\Mat_n(\R)$. Then:
\begin{enumerate}[label=(\roman*)]
  \item $\abs{\det(A)}$ equals the $n$-dimensional volume of the
    parallelepiped
    $P = \setcond{t_1 v_1+\cdots+t_n v_n}{0\leq t_i\leq 1}$.
  \item $\det(A)>0$ if and only if the ordered basis
    $(v_1,\ldots,v_n)$ has the same orientation as the standard basis.
\end{enumerate}
\end{theorem}

\begin{proof}[Proof sketch]
(i) Both $\abs{\det}$ and the volume function are multiplicative
under composition ($\abs{\det(AB)}=\abs{\det A}\abs{\det B}$), agree
on elementary matrices (which correspond to elementary row
operations with known geometric effects), and every invertible matrix
is a product of elementary matrices. For the singular case, both
sides equal zero.

(ii) The determinant is a continuous function of the columns.
The set of invertible matrices $\GL_n(\R)$ has exactly two connected
components, distinguished by the sign of $\det$. The identity matrix
$I_n$ has $\det=1>0$ and corresponds to the standard orientation.
\end{proof}

\begin{center}
\begin{tikzpicture}[>=Stealth, scale=1.0]
  % 2D parallelogram
  \begin{scope}
    \coordinate (O) at (0,0);
    \coordinate (U) at (2.5,0.5);
    \coordinate (V) at (0.8,2);
    \coordinate (UV) at ($(U)+(V)$);

    \fill[blue!15] (O) -- (U) -- (UV) -- (V) -- cycle;
    \draw[thick, blue!70!black, ->] (O) -- (U)
      node[below right] {$\vec{u}$};
    \draw[thick, red!70!black, ->] (O) -- (V)
      node[above left] {$\vec{v}$};
    \draw[dashed, gray] (U) -- (UV) -- (V);
    \node at (1.6,1.2) {Area $= \abs{\det}$};
    \node[below] at (1.6,-0.5) {\textbf{2D:} parallelogram};
  \end{scope}

  % 3D parallelepiped
  \begin{scope}[shift={(7,0)}, x={(1,0)}, y={(0.4,0.5)}, z={(0,1)}]
    \coordinate (O) at (0,0,0);
    \coordinate (A) at (2,0,0);
    \coordinate (B) at (0,1.5,0);
    \coordinate (C) at (0,0,1.8);
    \coordinate (AB) at (2,1.5,0);
    \coordinate (AC) at (2,0,1.8);
    \coordinate (BC) at (0,1.5,1.8);
    \coordinate (ABC) at (2,1.5,1.8);

    % Back faces
    \fill[red!8] (O) -- (B) -- (AB) -- (A) -- cycle;
    \fill[blue!8] (O) -- (B) -- (BC) -- (C) -- cycle;
    \fill[green!8] (O) -- (A) -- (AC) -- (C) -- cycle;

    % Front faces
    \fill[red!15] (C) -- (BC) -- (ABC) -- (AC) -- cycle;
    \fill[blue!15] (A) -- (AB) -- (ABC) -- (AC) -- cycle;
    \fill[green!15] (B) -- (AB) -- (ABC) -- (BC) -- cycle;

    % Edges
    \draw[thick] (O) -- (A) -- (AB) -- (B) -- cycle;
    \draw[thick] (O) -- (C) -- (AC) -- (A);
    \draw[thick] (C) -- (BC) -- (B);
    \draw[thick] (AC) -- (ABC) -- (AB);
    \draw[thick] (ABC) -- (BC);

    % Vectors
    \draw[->, very thick, blue!70!black] (O) -- (A)
      node[below right] {$\vec{u}$};
    \draw[->, very thick, red!70!black] (O) -- (B)
      node[left] {$\vec{v}$};
    \draw[->, very thick, green!60!black] (O) -- (C)
      node[left] {$\vec{w}$};

    \node at (1.2,0.7,0.9) {\small Vol};
    \node[below] at (1,-0.2,-0.5) {\textbf{3D:} parallelepiped};
  \end{scope}
\end{tikzpicture}
\captionof{figure}{The determinant computes signed areas and volumes.}
\label{fig:det-area-volume}
\end{center}

\begin{center}
\begin{tikzpicture}[>=Stealth, scale=1.2]
  % Positive orientation
  \begin{scope}
    \draw[->] (-0.3,0) -- (2.5,0) node[right] {$x$};
    \draw[->] (0,-0.3) -- (0,2.2) node[above] {$y$};
    \draw[->, very thick, blue!70!black] (0,0) -- (2,0.3)
      node[below right] {$\vec{u}$};
    \draw[->, very thick, red!70!black] (0,0) -- (-0.2,1.8)
      node[above left] {$\vec{v}$};
    \draw[->, thick, green!50!black] (0.7,0.1) arc[start angle=8,
      end angle=84, radius=0.7];
    \node[green!50!black] at (0.7,0.9) {\small $+$};
    \node[below] at (1,-0.6) {$\det > 0$};
    \node[below] at (1,-1.0) {(positive orientation)};
  \end{scope}

  % Negative orientation
  \begin{scope}[shift={(5.5,0)}]
    \draw[->] (-0.3,0) -- (2.5,0) node[right] {$x$};
    \draw[->] (0,-0.3) -- (0,2.2) node[above] {$y$};
    \draw[->, very thick, blue!70!black] (0,0) -- (-0.2,1.8)
      node[above left] {$\vec{u}$};
    \draw[->, very thick, red!70!black] (0,0) -- (2,0.3)
      node[below right] {$\vec{v}$};
    \draw[->, thick, orange!70!red] (0.05,0.65) arc[start angle=84,
      end angle=8, radius=0.7];
    \node[orange!70!red] at (0.7,0.9) {\small $-$};
    \node[below] at (1,-0.6) {$\det < 0$};
    \node[below] at (1,-1.0) {(negative orientation)};
  \end{scope}
\end{tikzpicture}
\captionof{figure}{The sign of the determinant encodes orientation.}
\label{fig:det-orientation}
\end{center}


% ============================================================
\section{Applications}\label{sec:det-applications}
% ============================================================

\subsection{Cross product in $\R^3$}\label{subsec:cross-product}

\begin{definition}{Cross product}{def:cross-product}
\index{cross product}
For $\vec{u}=(u_1,u_2,u_3)$ and $\vec{v}=(v_1,v_2,v_3)$ in~$\R^3$,
the \emph{cross product} is defined (formally) by
\[
  \vec{u}\times\vec{v} \deff
  \det\begin{pmatrix}
    \vece{1} & \vece{2} & \vece{3} \\
    u_1 & u_2 & u_3 \\
    v_1 & v_2 & v_3
  \end{pmatrix}
  = \begin{pmatrix}
    u_2 v_3 - u_3 v_2 \\
    u_3 v_1 - u_1 v_3 \\
    u_1 v_2 - u_2 v_1
  \end{pmatrix}.
\]
\end{definition}

\begin{proposition}{Properties of the cross product}{prop:cross-product}
For $\vec{u},\vec{v},\vec{w}\in\R^3$ and $\alpha\in\R$:
\begin{enumerate}[label=(\roman*)]
  \item \textbf{Antisymmetry:}
    $\vec{v}\times\vec{u} = -(\vec{u}\times\vec{v})$.
  \item \textbf{Bilinearity:}
    $(\alpha\vec{u}+\vec{w})\times\vec{v}
    = \alpha(\vec{u}\times\vec{v})+\vec{w}\times\vec{v}$.
  \item \textbf{Orthogonality:}
    $(\vec{u}\times\vec{v})\cdot\vec{u}
    = (\vec{u}\times\vec{v})\cdot\vec{v} = 0$.
  \item \textbf{Norm:}
    $\norm{\vec{u}\times\vec{v}}$ equals the area of the
    parallelogram Spanned by $\vec{u}$ and~$\vec{v}$.
  \item \textbf{Scalar triple product:}
    $(\vec{u}\times\vec{v})\cdot\vec{w}
    = \det(\vec{u},\vec{v},\vec{w})$ (columns).
\end{enumerate}
\end{proposition}

\begin{proof}
(i)--(ii) follow from properties of the determinant.
(iii) If we substitute $\vec{u}$ for $\vec{w}$ in the triple product
$\det(\vec{u},\vec{v},\vec{w})$, the matrix has two identical columns,
so $\det=0$.
(iv) follows from the identity
$\norm{\vec{u}\times\vec{v}}^2 = \norm{\vec{u}}^2\norm{\vec{v}}^2
- (\vec{u}\cdot\vec{v})^2$ (Lagrange identity).
(v) is verified by direct computation: both sides equal
$u_1(v_2 w_3-v_3 w_2) - u_2(v_1 w_3-v_3 w_1) + u_3(v_1 w_2-v_2 w_1)$.
\end{proof}

\subsection{Area and volume formulas}\label{subsec:area-volume}

\begin{proposition}{Area of a triangle}{prop:triangle-area}
\index{area!triangle}
The area of a triangle with vertices $P_1=(x_1,y_1)$,
$P_2=(x_2,y_2)$, $P_3=(x_3,y_3)$ in $\R^2$ is
\[
  \text{Area} = \frac{1}{2}\abs*{\det\begin{pmatrix}
    x_2-x_1 & x_3-x_1 \\
    y_2-y_1 & y_3-y_1
  \end{pmatrix}}
  = \frac{1}{2}\abs*{\det\begin{pmatrix}
    x_1 & y_1 & 1 \\
    x_2 & y_2 & 1 \\
    x_3 & y_3 & 1
  \end{pmatrix}}.
\]
\end{proposition}

\begin{proof}
The triangle is half the parallelogram Spanned by
$\overrightarrow{P_1 P_2}$ and $\overrightarrow{P_1 P_3}$.
The second formula follows by cofactor expansion along the third
column of the $3\times 3$ matrix.
\end{proof}

\subsection{Vandermonde determinant}\label{subsec:vandermonde}

\begin{theorem}{Vandermonde determinant}{thm:vandermonde}
\index{Vandermonde determinant}
For $x_1,\ldots,x_n\in\K$,
\[
  V(x_1,\ldots,x_n) \deff
  \det\begin{pmatrix}
    1 & x_1 & x_1^2 & \cdots & x_1^{n-1} \\
    1 & x_2 & x_2^2 & \cdots & x_2^{n-1} \\
    \vdots & \vdots & \vdots & \ddots & \vdots \\
    1 & x_n & x_n^2 & \cdots & x_n^{n-1}
  \end{pmatrix}
  = \prod_{1\leq i < j \leq n} (x_j - x_i).
\]
In particular, $V(x_1,\ldots,x_n)\neq 0$ if and only if the
$x_i$ are pairwise distinct.
\end{theorem}

\begin{proof}
We proceed by induction on~$n$. For $n=1$, $V(x_1)=1$ and the empty
product is~$1$. For $n=2$, $V(x_1,x_2) = x_2 - x_1$.

For the inductive step, subtract $x_1$ times column $j-1$ from
column~$j$, for $j=n,n-1,\ldots,2$ (right to left). This does not
change the determinant. Row~$1$ becomes $(1,0,0,\ldots,0)$, and for
$i\geq 2$, the entry in row~$i$, column~$j$ becomes
$x_i^{j-1}-x_1 x_i^{j-2} = x_i^{j-2}(x_i - x_1)$.
Expanding along row~$1$ and factoring out $(x_i-x_1)$ from row~$i$
(for $i=2,\ldots,n$):
\[
  V(x_1,\ldots,x_n) = \prod_{i=2}^n (x_i-x_1)\cdot V(x_2,\ldots,x_n).
\]
By the induction hypothesis,
$V(x_2,\ldots,x_n) = \prod_{2\leq i<j\leq n}(x_j-x_i)$, so
\[
  V(x_1,\ldots,x_n)
  = \prod_{i=2}^n (x_i-x_1)\cdot\prod_{2\leq i<j\leq n}(x_j-x_i)
  = \prod_{1\leq i<j\leq n}(x_j-x_i). \qedhere
\]
\end{proof}

\begin{example}{Vandermonde for $n=3$}{ex:vandermonde-3}
\[
  \det\begin{pmatrix}1 & a & a^2 \\ 1 & b & b^2 \\ 1 & c & c^2
  \end{pmatrix}
  = (b-a)(c-a)(c-b).
\]
\end{example}

\begin{remark}{Application of Vandermonde}{rem:vandermonde-app}
The Vandermonde determinant appears in polynomial interpolation:
the system expressing a polynomial of degree $\leq n-1$ through
$n$ points $(x_i,y_i)$ has a unique solution if and only if
the $x_i$ are distinct, since the coefficient matrix is the
Vandermonde matrix.
\end{remark}


% ============================================================
\section{Exercises}\label{sec:det-exercises}
% ============================================================

\begin{exercise}{Determinants of small matrices \basic}{exo:det-small}
Compute the following determinants:
\begin{enumerate}[label=(\alph*)]
  \item $\det\begin{pmatrix}3 & 1 \\ 5 & 2\end{pmatrix}$.
  \item $\det\begin{pmatrix}2 & -1 & 3 \\ 0 & 4 & 1 \\ 1 & 2 & -1\end{pmatrix}$.
  \item $\det\begin{pmatrix}1 & 2 & 3 & 4 \\ 0 & 1 & 2 & 3 \\
    0 & 0 & 1 & 2 \\ 0 & 0 & 0 & 1\end{pmatrix}$.
\end{enumerate}
\end{exercise}

\begin{exercise}{Row operations and determinant \basic}{exo:row-ops-det}
Let $A\in\Mat_3(\R)$ with $\det(A)=6$. Compute $\det(B)$ where $B$ is
obtained from $A$ by:
\begin{enumerate}[label=(\alph*)]
  \item swapping rows 1 and 3,
  \item multiplying row 2 by $-4$,
  \item adding 5 times row 1 to row 2,
  \item performing all three operations in sequence.
\end{enumerate}
\end{exercise}

\begin{exercise}{Determinant and scalar multiplication \basic}{exo:det-scalar}
\index{determinant!scalar multiplication}
Let $A\in\Mat_n(\K)$ with $\det(A)=d$. Express the following in terms
of $d$ and~$n$:
\begin{enumerate}[label=(\alph*)]
  \item $\det(3A)$.
  \item $\det(-A)$.
  \item $\det(\transpose{A})$.
  \item $\det(A^2)$.
\end{enumerate}
\end{exercise}

\begin{exercise}{Determinant of a rank-1 update \intermediate}{exo:rank-1-update}
\index{matrix determinant lemma}
Let $A\in\GL_n(\K)$ and let $u,v\in\K^n$ (column vectors). Prove
the \emph{matrix determinant lemma}:
\[
  \det(A + uv^{\mathsf{T}}) = (1 + v^{\mathsf{T}}\inv{A}u)\det(A).
\]
\textit{Hint:} Factor $A+uv^{\mathsf{T}} = A(I+\inv{A}uv^{\mathsf{T}})$
and compute the determinant of $I+wv^{\mathsf{T}}$ for a column
vector $w=\inv{A}u$.
\end{exercise}

\begin{exercise}{Cofactor expansion practice \basic}{exo:cofactor-practice}
Compute the determinant of
\[
  A = \begin{pmatrix}
    2 & 0 & 1 & 3 \\
    1 & -1 & 0 & 2 \\
    3 & 2 & 0 & 1 \\
    0 & 4 & -1 & 0
  \end{pmatrix}
\]
by expanding along a row or column of your choice (choose wisely to
minimise work).
\end{exercise}

\begin{exercise}{Adjugate of a $2\times 2$ matrix \basic}{exo:adjugate-2x2}
For $A=\begin{psmallmatrix}a&b\\c&d\end{psmallmatrix}$, compute
$\operatorname{adj}(A)$ and verify that
$A\cdot\operatorname{adj}(A)=\det(A)\,I_2$.
\end{exercise}

\begin{exercise}{Cramer's rule application \intermediate}{exo:cramer-app}
Use Cramer's rule to solve the system
\[
  \begin{cases}
    x + 2y - z = 3,\\
    2x - y + 3z = 1,\\
    3x + y + 2z = 7.
  \end{cases}
\]
\end{exercise}

\begin{exercise}{Vandermonde determinant \intermediate}{exo:vandermonde}
\begin{enumerate}[label=(\alph*)]
  \item Compute the Vandermonde determinant $V(1,2,3,4)$.
  \item For which values of $a\in\R$ is the Vandermonde matrix
    with nodes $1,a,a^2$ singular?
\end{enumerate}
\end{exercise}

\begin{exercise}{Determinant and eigenvalues \intermediate}{exo:det-eigenvalues}
\index{eigenvalue!and determinant}
Let $A\in\Mat_n(\K)$.
\begin{enumerate}[label=(\alph*)]
  \item Show that $\lambda$ is an eigenvalue of $A$ if and only if
    $\det(A-\lambda I_n)=0$.
  \item If $A$ has eigenvalues $\lambda_1,\ldots,\lambda_n$
    (counted with algebraic multiplicity), prove that
    $\det(A)=\lambda_1\cdots\lambda_n$.
\end{enumerate}
\end{exercise}

\begin{exercise}{Determinant of a block diagonal matrix \basic}{exo:block-diag}
Let $A=\diag(A_1,A_2,\ldots,A_k)$ be a block diagonal matrix where
$A_i\in\Mat_{n_i}(\K)$. Prove that
$\det(A)=\det(A_1)\det(A_2)\cdots\det(A_k)$.
\end{exercise}

\begin{exercise}{Antisymmetric matrices \intermediate}{exo:antisymmetric}
\index{antisymmetric matrix!determinant}
Let $A\in\Mat_n(\R)$ be \emph{antisymmetric}, i.e.\
$\transpose{A}=-A$.
\begin{enumerate}[label=(\alph*)]
  \item Show that if $n$ is odd, then $\det(A)=0$.
  \item Give an example of a $2\times 2$ antisymmetric matrix
    with $\det(A)\neq 0$.
  \item Show that for even~$n$, $\det(A)\geq 0$.
    \textit{Hint:} Consider $\det(A)=\det(\transpose{A})=\det(-A)$.
\end{enumerate}
\end{exercise}

\begin{exercise}{Derivative of the determinant \challenging}{exo:det-derivative}
\index{determinant!derivative}
Let $A(t)=(a_{ij}(t))$ be an $n\times n$ matrix whose entries are
differentiable functions of $t\in\R$. Prove \emph{Jacobi's formula}:
\[
  \frac{d}{dt}\det(A(t))
  = \sum_{i=1}^n \det\bigl(A_1(t),\ldots,A_i'(t),\ldots,A_n(t)\bigr),
\]
where $A_i(t)$ denotes the $i$-th row of $A(t)$ and $A_i'(t)$ its
derivative.
\textit{Hint:} Use multilinearity of $\det$ in the rows and the
product rule.
\end{exercise}

\begin{exercise}{Cauchy determinant \challenging}{exo:cauchy-det}
\index{Cauchy determinant}
Let $x_1,\ldots,x_n$ and $y_1,\ldots,y_n$ be elements of a
field~$\K$ with $x_i+y_j\neq 0$ for all $i,j$. The \emph{Cauchy
matrix} is $C=(c_{ij})$ with $c_{ij}=\frac{1}{x_i+y_j}$.
Prove the \emph{Cauchy determinant formula}:
\[
  \det(C) = \frac{\prod_{1\leq i<j\leq n}(x_j-x_i)(y_j-y_i)}%
  {\prod_{i,j=1}^{n}(x_i+y_j)}.
\]
\textit{Hint:} Proceed by induction, reducing to a Vandermonde-type
computation.
\end{exercise}

\begin{exercise}{Determinant of the commutator \challenging}{exo:commutator-det}
Let $A,B\in\Mat_n(\K)$.
\begin{enumerate}[label=(\alph*)]
  \item Show that if $n=2$ and $\tr(A)=\tr(B)=0$, then
    $\det(AB-BA)\leq 0$ when $\K=\R$.
  \item Construct $A,B\in\Mat_3(\R)$ such that $\det(AB-BA)\neq 0$.
  \item Show that for any $n$, $\tr(AB-BA)=0$.
\end{enumerate}
\end{exercise}

\begin{exercise}{Polynomial identity for determinants \intermediate}{exo:det-poly}
Let $A,B\in\Mat_n(\K)$. Show that
\[
  \det(A+B) + \det(A-B) = 2\sum_{k\text{ even}}D_k(A,B),
\]
where $D_k(A,B)$ is the sum of all determinants obtained by choosing
$k$ rows from $B$ and the remaining $n-k$ rows from $A$ (in their
original positions). Verify this for $n=2$ explicitly.
\end{exercise}


% ============================================================
\section*{Chapter summary}\label{sec:summary-ch4}
\addcontentsline{toc}{section}{Chapter summary}
% ============================================================

\begin{itemize}
  \item A \textbf{permutation}\index{permutation} of $\set{1,\ldots,n}$
    is a bijection; the set of all permutations forms the
    \textbf{symmetric group}~$\mathfrak{S}_n$ of order $n!$.
    Every permutation has a well-defined
    \textbf{signature}\index{signature!of a permutation}
    $\sgn(\sigma)=\pm 1$ (
ef{def:def:signature}).

  \item The \textbf{determinant}\index{determinant} of
    $A\in\Mat_n(\K)$ is defined by the \textbf{Leibniz formula}
    (
ef{def:def:determinant}):
    $\det(A) = \sum_{\sigma\in\mathfrak{S}_n}\sgn(\sigma)
    \prod_{i=1}^n a_{i,\sigma(i)}$.

  \item The determinant is the \emph{unique} function that is
    \textbf{multilinear}, \textbf{alternating} in the rows, and
    satisfies $\det(I_n)=1$ (
ef{thm:thm:det-unique}).

  \item Key \textbf{properties}: $\det(\transpose{A})=\det(A)$
    (
ef{thm:thm:det-transpose}), $\det(AB)=\det(A)\det(B)$
    (
ef{thm:thm:det-product}), triangular matrices have
    $\det=\prod a_{ii}$ (
ef{prop:prop:det-triangular}).

  \item \textbf{Cofactor (Laplace) expansion}\index{Laplace expansion}
    computes $\det$ by expanding along any row or column
    (
ef{thm:thm:laplace}).

  \item The \textbf{adjugate matrix}\index{adjugate matrix} satisfies
    $A\cdot\operatorname{adj}(A) = \det(A)\,I_n$
    (
ef{thm:thm:adjugate}), yielding the formula
    $\inv{A}=\frac{1}{\det(A)}\operatorname{adj}(A)$.

  \item \textbf{Cramer's rule}\index{Cramer's rule} gives an explicit
    formula for the solution of $Ax=b$ when $\det(A)\neq 0$
    (
ef{thm:thm:cramer}).

  \item $A$ is \textbf{invertible} if and only if $\det(A)\neq 0$
    (
ef{thm:thm:det-invertibility}).

  \item The determinant of a \textbf{linear map} is well-defined
    (independent of the choice of basis, 
ef{prop:prop:det-well-defined}).

  \item \textbf{Geometrically}, $\abs{\det(A)}$ is the volume of the
    parallelepiped Spanned by the columns, and $\sgn(\det)$ encodes
    orientation (
ef{thm:thm:signed-volume}).

  \item \textbf{Applications}: the cross product in~$\R^3$, area and
    volume formulas, and the \textbf{Vandermonde
    determinant}\index{Vandermonde determinant}
    $\prod_{i<j}(x_j-x_i)$ (
ef{thm:thm:vandermonde}).
\end{itemize}
